%% ============================================================
%% Evaluation Section — Quantitative Results
%% "Trustless-by-Design Security Framework for Distributed IoT
%%  of 6G Non-Terrestrial Networks" — IEEE TNSM/IoTJ Extension
%% ============================================================

\section{Quantitative Performance Evaluation}
\label{sec:evaluation}

This section presents the quantitative simulation results for the six
experiment groups introduced in Section~\ref{sec:simulation}.  All
experiments run 30 independent Monte-Carlo trials (seed offset $s+r$,
$r\in\{0,\dots,29\}$) to report $95\%$ confidence intervals (CI).
The complete simulation codebase, configuration, and raw results are
available at~\cite{repo2026}.

% ----------------------------------------------------------
\subsection{Trust Dynamics and Detection Latency (Exp1)}
% ----------------------------------------------------------

Fig.~\ref{fig:trust_evolution}(a) plots the trust score assigned to
a representative malicious \emph{build-then-betray} node by each
model over 500 rounds.  The proposed framework decays the score to
below the detection threshold $\theta = 0.4$ within
\textbf{35.9\,$\pm$\,5.6 rounds} (mean\,$\pm$\,95\%\,CI) of betrayal
onset, compared to 85.5\,$\pm$\,7.8 rounds for ZT~Auth-Only,
152.6\,$\pm$\,12.8 rounds for Static~DLT, and 219.8\,$\pm$\,16.3
rounds for Centralized~Trust (Table~\ref{tab:comparison}).  The
improvement over Centralized~Trust is $6.1\times$ because the
authority is unreachable during partition and cannot revoke trust in
real time.  The Static~DLT improvement of $4.3\times$ arises from
the absence of temporal decay: once a node accumulates reputation, it
is not penalised for subsequent misbehaviour without explicit
revocation.

Fig.~\ref{fig:trust_evolution}(b) confirms that legitimate nodes
maintain stable trust trajectories during normal operation despite the
aggressive per-segment decay introduced by Eq.~(4).

% ----------------------------------------------------------
\subsection{Partition Resilience (Exp2)}
% ----------------------------------------------------------

We inject a 50-round network partition that isolates $30\%$ of nodes.
The proposed DLT reconciliation mechanism applies trust aging to
stale records upon reconnection (Algorithm~\ref{alg:reconcile}),
preventing malicious nodes from carrying over inflated pre-partition
scores.

Table~\ref{tab:partition} and Fig.~\ref{fig:partition}(a) show that
trust divergence (mean absolute deviation between partitioned and
reference sub-graphs) rises during the partition and collapses to
near zero within $\approx 5$ rounds post-reconnection.
Fig.~\ref{fig:partition}(b) shows that detection F1 \emph{improves}
by $+0.130$ after partition resolution because stale malicious trust
scores are penalised by accumulated time-decay, exposing attackers
that had built reputation before the partition.

% ----------------------------------------------------------
\subsection{Cryptographic Overhead (Exp3)}
% ----------------------------------------------------------

Table~\ref{tab:crypto} reports the Paillier~HE overhead for
privacy-preserving trust aggregation (Eq.~1--2).  With a
1024-bit key, aggregating 100 contributor ciphertexts costs
\textbf{1{,}499.6\,ms} total (encryption + aggregation + decryption)
and produces a 25.0\,KB aggregate ciphertext.  Overhead scales
\emph{linearly} with contributor count, consistent with the additive
homomorphic property.  For 50-node clusters (typical NTN cell),
latency is 755\,ms per aggregation cycle — well within the 60-second
round duration adopted in this study.

The ledger consensus throughput reaches \textbf{1{,}872\,TPS} at
batch size 100, decreasing to 327\,TPS at batch size 10 due to
per-transaction PBFT-like overhead.  All throughput figures exceed
the trust record generation rate of $N/60 \approx 17$\,records/s for
$N=1{,}000$ nodes.

% ----------------------------------------------------------
\subsection{Privacy-Utility Tradeoff (Exp4)}
% ----------------------------------------------------------

Fig.~\ref{fig:privacy} plots detection F1 and trust estimation MAE
as functions of the DP budget $\varepsilon$.  A tight privacy budget
($\varepsilon = 0.1$, $\sigma = 48.4$) introduces noise that raises
MAE to $0.499$ but F1 remains at $0.291\,\pm\,0.010$ — \emph{above}
the no-DP baseline of $0.249$.  This counter-intuitive improvement
occurs because high-noise trust scores shift the operating point
favourably for the mixed attack population; under strict privacy, the
Gaussian noise masks on-off attackers' temporary ``good'' phases,
causing earlier detection.  At the practical operating point
$\varepsilon = 1.0$ ($\sigma = 4.85$), F1 is $0.281$ with MAE $0.480$.

% ----------------------------------------------------------
\subsection{Threshold Calibration (Analytical Contribution)}
% ----------------------------------------------------------
\label{subsec:threshold}

A key insight from Exp1 is that the detection threshold $\theta$ must
be calibrated to the \emph{steady-state trust} of legitimate nodes,
which depends on the per-segment decay rate $\lambda_s$ and the
effective inter-update interval $\Delta t_{\mathrm{eff}} =
T_{\mathrm{round}} + \bar{d}_s/1000$.

Setting $T_{ij}(t) = T_{ij}(t+\Delta t)$ in Eq.~(4) for a
legitimate node with mean feedback score $\bar{s}$:
%
\begin{equation}
  T^*_s = \frac{(1-\alpha)\,\bar{s}}
               {1 - \alpha\,e^{-\lambda_s\,\Delta t_{\mathrm{eff}}}}
  \label{eq:steady_state}
\end{equation}
%
Using $\alpha=0.7$, $T_{\mathrm{round}}=60$\,s, $\bar{s}=0.85$, and
segment-specific $\lambda_s$ values from Table~\ref{tab:lambda}, we
obtain $T^*_{\mathrm{GEO}}=0.530$, $T^*_{\mathrm{Ground}}=0.356$,
$T^*_{\mathrm{LEO}}=0.323$, $T^*_{\mathrm{UAV}}=0.264$.  For
malicious nodes ($\bar{s}=0.10$), the corresponding floor values
range from $T^*_{\mathrm{mal}}=0.031$ (UAV) to $0.062$ (GEO).

The segment-optimal threshold that minimises both FPR and FNR is:
%
\begin{equation}
  \theta^*_s = \frac{T^*_{s,\mathrm{legit}} + T^*_{s,\mathrm{mal}}}{2}
  \label{eq:theta_opt}
\end{equation}
%
yielding $\theta^*_{\mathrm{Ground}}=0.199$,
$\theta^*_{\mathrm{GEO}}=0.296$.
A global threshold of $\theta=0.20$ satisfies
$T^*_{s,\mathrm{mal}} < 0.20 < T^*_{s,\mathrm{legit}}$ for all
segment types simultaneously, reducing the false positive rate from
$26.8\%$ (at $\theta=0.40$) to $\approx 0$ analytically.

\emph{In this study, $\theta=0.40$ was chosen to demonstrate the
conservative bound; adaptive per-segment threshold calibration via
Eq.~(\ref{eq:theta_opt}) is recommended in deployment and is
validated analytically here.}

% ----------------------------------------------------------
\subsection{Scalability (Exp5)}
% ----------------------------------------------------------

Fig.~\ref{fig:scalability}(a) shows computation time per round
growing from 5.1\,ms ($N=100$) to 227.7\,ms ($N=5{,}000$),
consistent with $\mathcal{O}(N)$ scaling (sparse NTN graph, mean
degree $\bar{d}\approx 4.7$).  Ledger size grows linearly from
0.95\,MB to 47.7\,MB over the same range (Fig.~\ref{fig:scalability}(b)).
At $N=5{,}000$, total memory overhead is $112$\,MB — well within
the resource budget of an edge computing node co-located with a
HAPS or LEO ground station.

% ----------------------------------------------------------
\subsection{Attack Resilience (Exp6)}
% ----------------------------------------------------------

Fig.~\ref{fig:attack}(a) plots detection F1 versus malicious
fraction for all five models across six attack densities (5--30\%).
The proposed framework and UAV+FL achieve comparable F1 (0.086--0.286)
and both improve monotonically with malicious fraction — a natural
effect of more malicious nodes generating more observable anomalies.
Centralized~Trust plateaus at F1\,$\approx$\,0.08 because partition
prevents real-time revocation; Static~DLT is bounded by reputation
inertia at F1\,$\approx$\,0.14.

Although UAV+FL achieves the lowest raw detection latency (17.3
rounds), it applies the aggressive UAV decay rate ($\lambda=0.05$)
uniformly across all segment types, causing GEO legitimate nodes to
be falsely penalised (FPR\,=\,0.279).  The proposed model applies
per-segment $\lambda_s$, correctly distinguishing fast-moving UAVs
from stable GEO satellites and achieving equal F1 with significantly
lower detection latency than all DLT-only baselines.

% ----------------------------------------------------------
% Tables
% ----------------------------------------------------------

\begin{table}[t]
\centering
\caption{Quantitative Comparison of Trust Models at 20\% Malicious
         Node Fraction ($N=1{,}000$ nodes, 30 runs, 95\% CI)}
\label{tab:comparison}
\begin{tabular}{lrrrr}
\toprule
Model & Det.\ Lat.\ (rds) & F1 & AUC & FPR \\
\midrule
\textbf{Proposed (DLT+ZTA)} & \textbf{35.9\,$\pm$\,5.6}  & 0.236 & 0.520 & 0.268 \\
ZT Auth-Only                 & 85.5\,$\pm$\,7.8            & 0.251 & 0.571 & 0.000 \\
Static DLT (NxGenT)          & 152.6\,$\pm$\,12.8          & 0.159 & 0.630 & 0.000 \\
Centralized Trust            & 219.8\,$\pm$\,16.3          & 0.091 & 0.623 & 0.000 \\
UAV+FL Blockchain            & 17.3\,$\pm$\,4.0            & 0.236 & 0.514 & 0.279 \\
\bottomrule
\end{tabular}
\end{table}

\begin{table}[t]
\centering
\caption{Paillier Homomorphic Encryption Overhead (1024-bit key)}
\label{tab:crypto}
\begin{tabular}{rrr}
\toprule
Contributors & Total Time (ms) & Ciphertext Size (KB) \\
\midrule
  5  &   81.1 &  1.25 \\
 10  &  155.3 &  2.50 \\
 25  &  385.2 &  6.25 \\
 50  &  755.3 & 12.50 \\
100  & 1{,}499.6 & 25.00 \\
200  & 3{,}037.1 & 50.00 \\
\bottomrule
\end{tabular}
\end{table}

\begin{table}[t]
\centering
\caption{Segment-Specific Decay Parameters and Calibrated Thresholds}
\label{tab:lambda}
\begin{tabular}{lrrrr}
\toprule
Segment & $\lambda_s$ (s$^{-1}$) & $T^*_{\mathrm{legit}}$ &
          $T^*_{\mathrm{mal}}$ & $\theta^*_s$ \\
\midrule
GEO    & 0.005 & 0.530 & 0.062 & 0.296 \\
Ground & 0.015 & 0.356 & 0.042 & 0.199 \\
LEO    & 0.020 & 0.323 & 0.038 & 0.181 \\
UAV    & 0.050 & 0.264 & 0.031 & 0.148 \\
\midrule
\multicolumn{4}{l}{Global ($\theta=0.20$: valid for all segments)} & 0.200 \\
\bottomrule
\end{tabular}
\end{table}
